\chapter{Threats to Validity}
\label{ch:threats-to-validity}

Any study based on publicly available artifacts is constrained by what those artifacts can reveal. The threats below are organized by threat type: reliability, construct validity, internal validity, and external validity.

\section{Reliability}
\label{sec:threat-reliability}

All measurements are based on public evidence collected in May 2026 and reflect a snapshot of each repository at that moment (Section~\ref{sec:open-source-repository-measurement}); the same project measured earlier or later might score differently. Some quality questions also require judgement, such as whether documentation is adequate or project metadata complete enough to support a score, and these were performed by a single assessor without an independent inter-rater check. The rationale is recorded in the measurement template (Section~\ref{subsec:measurement-template}), but the impression scores summarize observations beyond the recorded checklist answers, so the recorded data alone cannot fully reconstruct them. The time budget of one to four hours per package bounds the depth of every check; the surface measures are shallow by design.

\section{Construct Validity}
\label{sec:threat-construct-validity}

The measurements use public artifacts as evidence for software qualities, a practical proxy that does not capture every aspect of the underlying constructs: visible testing evidence does not prove that a package is correct, and visible issue-tracking activity does not fully describe maintainability. Mention counts likewise measure how often a package appears across candidate sources, not how good the software is (Section~\ref{sec:mention-count-ranking}).

\section{Internal Validity}
\label{sec:threat-internal-validity}

The candidate set was built from academic survey papers, community-curated lists, and LLM-assisted discovery, with the LLM lists supplementing rather than driving the other two sources. Those lists reflect the training data of the model rather than an independent survey of the MPS domain, so a package unknown to the model could be absent from its output. The filtering decisions described in Section~\ref{sec:filtering-criteria} further shape the measured set, since a package can be excluded for being out of scope, closed source, inactive, or not measurable from public repositories.

\section{External Validity}
\label{sec:threat-external-validity}

The selected set covers open-source MPS software measurable through public artifacts. Proprietary platforms, discontinued projects, and narrowly scoped libraries may differ from the measured packages, so the results describe the selected open-source MPS sample, not all robotics software. Because the packages play different roles, comparisons are evidence about development practice across a domain rather than direct product-to-product comparisons.

The installation-based measurements were performed on a Linux virtual machine, so the installability, surface reliability, and surface robustness results reflect Linux behaviour only. The same packages may install and behave differently on Windows or macOS, and installation routes available at measurement time, such as distribution packages, may be absent on other operating systems or later distribution releases.
